% Adapted from the author-supplied normalized-valuation-links.pdf,
% 23 September 2026. Source bytes and review scope are recorded separately.
\subsection{Excellent germs and boundary ideals}
\label{sec:local-excellent-boundary-links}

The normalized boundary itself has a construction in every dimension.
In this subsection let $(R,\mathfrak m)$ be a positive-dimensional
excellent reduced Noetherian local ring. A nonnegative real semivaluation
has $v(0)=+\infty$, $v(1)=0$, $v(fg)=v(f)+v(g)$ and
$v(f+g)\geq\min\{v(f),v(g)\}$. We use the centered cone
\[
 \mathcal C_R
   :=\{v:R\longrightarrow[0,+\infty]\mid
          v\text{ is a semivaluation},\ 0<v(\mathfrak m)<+\infty\}.
\]
It has the topology of pointwise convergence, with the compact order
topology on $[0,+\infty]$, and the positive scaling action $v\mapsto tv$.
For an ideal $I=(f_1,\ldots,f_N)$, set $v(I)=\min_jv(f_j)$.
Expressing either generating set in terms of the other proves independence
of this choice. Units have value zero. The support
$\{f:v(f)=+\infty\}$ is prime; it need not be a minimal prime.
For an $\mathfrak m$-primary ideal $\mathfrak a$, write
$\mathcal V_{\mathfrak a}(R)=\{v\in\mathcal C_R:v(\mathfrak a)=1\}$,
consistently with the curve notation below.

\begin{theorem}\label{thm:p121-excellent-primary-links}
Every $\mathcal V_{\mathfrak a}(R)$ is nonempty compact Hausdorff,
and $(t,v)\mapsto tv$ is a homeomorphism
$(0,+\infty)\times\mathcal V_{\mathfrak a}(R)\simeq\mathcal C_R$.
For primary ideals $\mathfrak a,\mathfrak b$, the unique map of
sections preserving every ray is the homeomorphism
\[
 H_{\mathfrak a,\mathfrak b}(v)=\frac{v}{v(\mathfrak b)}.
\tag{70.7a}\label{eq:p121-general-boundary-change}
\]
It has inverse $H_{\mathfrak b,\mathfrak a}$ and satisfies
$H_{\mathfrak b,\mathfrak d}H_{\mathfrak a,\mathfrak b}
=H_{\mathfrak a,\mathfrak d}$ for every third primary ideal.
Thus the quotient $\mathcal C_R/\mathbf R_{>0}$ and these sections
are independent of generators and birational models.
\end{theorem}

\begin{proof}
The semivaluation axioms define a closed subset of
$[0,+\infty]^R$: extended addition and minimum are continuous,
and the order relation is closed. The additional condition
$v(\mathfrak m)=1$ is closed because $\mathfrak m$ is finitely
generated. Hence $\mathcal V_{\mathfrak m}(R)$ is compact Hausdorff.
Choose $r\geq s\geq1$ with
$\mathfrak m^r\subseteq\mathfrak a\subseteq\mathfrak m^s$.
For every centered $v$,
\[
                  s\,v(\mathfrak m)\leq v(\mathfrak a)
                                     \leq r\,v(\mathfrak m).
\tag{70.7b}\label{eq:p121-primary-gauge-bounds}
\]
Thus $v(\mathfrak a)$ is positive and finite. The continuous map
$v\mapsto v/v(\mathfrak a)$ sends
$\mathcal V_{\mathfrak m}(R)$ onto $\mathcal V_{\mathfrak a}(R)$,
proving compactness. The inverse of the asserted cone homeomorphism
is $u\mapsto(u(\mathfrak a),u/u(\mathfrak a))$; positive scaling
is jointly continuous also at infinite coordinate values.

If $\mathfrak m^u\subseteq\mathfrak b\subseteq\mathfrak m^t$,
then $t/r\leq v(\mathfrak b)\leq u/s$ on
$\mathcal V_{\mathfrak a}(R)$. This proves continuity of the change
of section with a denominator bounded away from zero and infinity.
Homogeneity proves its inverse and cocycle formulas. On the ray of $v$,
the only point having $\mathfrak b$-value one is $v/v(\mathfrak b)$.
The cone product decomposition identifies its ray quotient with any section.

For nonemptiness, choose a saturated prime chain ending at $\mathfrak m$
and let $\mathfrak p$ be its penultimate member. Then $R/\mathfrak p$
is an excellent one-dimensional local domain. Excellence implies the
Nagata property, and normalization is finite
\cite[Tags 07QV and 035S]{StacksProject}. Localizing that normalization
at a maximal ideal above $\mathfrak m/\mathfrak p$ gives a
one-dimensional normal Noetherian local domain, hence a DVR
\cite[Tag 00PD]{StacksProject}. Its order pulls back to an element of
$\mathcal C_R$, which can be rescaled into every primary section.
\end{proof}

Let $S=\prod_iS_i$ be the normalization of $R$ in its total ring of
fractions. It is finite, and $S_i$ is the normal domain associated to a
minimal component; see \cite[Tags 035P and 035S]{StacksProject}.
The conductor $\mathfrak c=\operatorname{Ann}_R(S/R)$ is also an
ideal of $S$, with $\mathfrak cS=\mathfrak c$. In particular
\[
                    R=S\mathbin{\times}_{S/\mathfrak cS}R/\mathfrak c.
\tag{70.7c}\label{eq:p121-conductor-pullback-square}
\]
Indeed, an element of $S$ whose residue comes from $R/\mathfrak c$
differs from an element of $R$ by an element of $\mathfrak cS\subset R$.

Let $\mathcal B_{\mathfrak a,i}$ be the nonnegative real semivaluations
on $S_i$ with $w(\mathfrak aS_i)=1$, and set
$\mathcal B_{\mathfrak a}=\coprod_i\mathcal B_{\mathfrak a,i}$.
Each summand is closed in $[0,+\infty]^{S_i}$, so this finite disjoint
union is compact Hausdorff. Restriction defines a continuous map
$\pi_{\mathfrak a}:\mathcal B_{\mathfrak a}\to
\mathcal V_{\mathfrak a}(R)$; the primary-ideal bounds prove the
centering condition downstairs.

This presentation retains all local normalization branches. The center
$\{s:w(s)>0\}$ contains $\mathfrak mS_i$ and is a maximal ideal
$\mathfrak n$ above $\mathfrak m$, since $S_i/\mathfrak mS_i$ is
finite over $R/\mathfrak m$. Every denominator outside $\mathfrak n$
has value zero, so $w$ extends uniquely to $(S_i)_{\mathfrak n}$.
There are finitely many such centers. Their pieces are clopen:
the Chinese remainder theorem supplies an element equal to zero modulo
one of these maximal ideals and to one modulo all the others; its
positive-value locus selects that center. Restriction from the local
branch link is a homeomorphism onto this piece, since localized
evaluations are the corresponding numerator evaluations.

\begin{lemma}\label{lem:p121-normalization-semivaluation-lift}
The restriction map $\pi_{\mathfrak a}$ is surjective.
\end{lemma}

\begin{proof}
Let $v$ have support $\mathfrak p$. Choose a minimal prime
$\mathfrak q_i\subseteq\mathfrak p$ and, by lying over, a prime
$P\subset S_i$ above $\mathfrak p/\mathfrak q_i$.
The field $L=\operatorname{Frac}(S_i/P)$ is finite over
$K=\operatorname{Frac}(R/\mathfrak p)$. The finite-valued part of
$v$ defines a valuation ring $A\subset K$ with value group
$\Gamma_A\subset\mathbf R$.

Let $A'$ be the integral closure of $A$ in $L$, choose a prime
$\mathfrak n\subset A'$ over the maximal ideal of $A$, and take
a valuation ring $W\subset L$ dominating $A'_{\mathfrak n}$
\cite[Tag 00IA]{StacksProject}. Then $W\cap K$ dominates $A$
and equals $A$: if $x\in(W\cap K)\setminus A$, its inverse
lies in the maximal ideal of $A$, hence in that of $W$, contradicting
$x\in W$.

The inclusion of value groups has finite index
\cite[Tag 0ASH]{StacksProject}. For $\gamma\in\Gamma_W$, choose
$N>0$ with $N\gamma\in\Gamma_A$ and assign to $\gamma$ the real
number $(N\gamma)/N$, using the given embedding of $\Gamma_A$
in $\mathbf R$. A common multiple of two choices proves independence;
the same construction respects addition and order. It therefore extends
the original real valuation. Every element of $S_i/P$ is integral
over $R/\mathfrak p\subset A$ and belongs to the integrally closed
ring $W$. Restrict this real valuation to $S_i/P$ and give elements
of $P$ value $+\infty$. The resulting semivaluation restricts to
$v$ and has $w(\mathfrak aS_i)=1$.
\end{proof}

\begin{theorem}\label{thm:p121-excellent-conductor-descent}
The map $\pi_{\mathfrak a}$ presents $\mathcal V_{\mathfrak a}(R)$
as the compact Hausdorff quotient of $\mathcal B_{\mathfrak a}$ by
the closed relation of equal restrictions to $R$.
Every fiber over $v$ with $v(\mathfrak c)<+\infty$ is a singleton.
Put
\[
 Z_R=\{v\in\mathcal V_{\mathfrak a}(R):v(\mathfrak c)=+\infty\},
 \qquad Z_B=\pi_{\mathfrak a}^{-1}(Z_R).
\]
These are closed. On $Z_B$, all lifts factor through $S/\mathfrak cS$,
and equality of their restrictions to $R/\mathfrak c$ is exactly the
remaining equivalence relation. There is a topological pushout
\[
              \mathcal V_{\mathfrak a}(R)
                  \simeq\mathcal B_{\mathfrak a}\sqcup_{Z_B}Z_R.
\]
Restriction, the conductor strata, and this quotient presentation commute
with every primary change \eqref{eq:p121-general-boundary-change}.
\end{theorem}

\begin{proof}
The continuous surjection in
\cref{lem:p121-normalization-semivaluation-lift} is a quotient map
because its domain is compact and its target is Hausdorff.
Its equivalence relation is the inverse image of the diagonal, hence closed.
Regard a point of $\mathcal B_{\mathfrak a,i}$ as a semivaluation on
$S$ by projection to $S_i$. This loses no information: multiplicativity
and $1=\sum_ie_i$ force a semivaluation on a finite product to take
value zero on exactly one primitive idempotent $e_i$ and infinity on
all the others.

If $v(\mathfrak c)<+\infty$, choose $c\in\mathfrak c$ with
$v(c)<+\infty$. Since $cs\in R$ for every $s\in S$, any lift satisfies
\[
                         w(s)=v(cs)-v(c).
\]
There is no subtraction of infinite quantities here. This determines
both the component and the lift uniquely. If $v(\mathfrak c)=+\infty$,
both $v$ and its lifts annihilate the conductor. Their equality on $R$
is precisely equality after restriction to $R/\mathfrak c$, consistently
with \eqref{eq:p121-conductor-pullback-square}. Finite generation of
$\mathfrak c$ makes the infinite-value condition closed. The indicated
pushout is a quotient of the compact space
$\mathcal B_{\mathfrak a}\coprod Z_R$, and its induced map to
$\mathcal V_{\mathfrak a}(R)$ is a continuous bijection. Compactness
and Hausdorffness make it a homeomorphism.
Finally $w(\mathfrak bS_i)=v(\mathfrak b)$, so all changes of section
commute with restriction and preserve the conductor strata.
\end{proof}

\begin{example}\label{ex:p121-surface-conductor-fiber}
For $R=k[[x,y,z]]/(xy)$, normalization is
$k[[x,z]]\times k[[y,z]]$. The image consists of pairs agreeing
on the $z$-axis, and the conductor is $(x,y)$: multiplying a pair
by the two primitive idempotents tests that its common $z$-axis
restriction vanishes. The semivaluation obtained from
$R/(x,y)=k[[z]]$ by $\ord_z$ has $v(\mathfrak m)=1$ and
$v(\mathfrak c)=+\infty$. It has two distinct lifts, with
$w(x)=+\infty$ on the first branch and $w(y)=+\infty$ on the
second. Thus conductor identifications really occur in dimension two.
The theorem confines identifications to the conductor stratum; it does
not assert that every point of that stratum has more than one lift.
\end{example}

\begin{proposition}\label{prop:p121-blowup-boundary-gauge}
Let $I=(f_1,\ldots,f_N)$ be an ideal of $R$. If $v\in\mathcal C_R$
and $v(I)<+\infty$, then $v$ has a canonical center on
$\operatorname{Bl}_I(\operatorname{Spec}R)$, and its value on the
invertible ideal $I\mathcal O_{\operatorname{Bl}_I R}$ is $v(I)$.
\end{proposition}

\begin{proof}
Choose $j$ with $v(f_j)=v(I)<+\infty$. The relevant blowup chart is
\[
                  R[f_1/f_j,\ldots,f_N/f_j]\subset R[1/f_j].
\]
The valuation ring of $\operatorname{Frac}(R/\operatorname{supp}v)$
receives this chart, since $v(f_i/f_j)=v(f_i)-v(f_j)\geq0$.
It determines a center there. Two minimizing generators have a ratio
of value zero, so the centers agree on their chart overlap. A generator
with larger finite value gives a negative-valued ratio and cannot give
a center; a generator of infinite value vanishes at the generic
support point and its chart cannot receive that point. On a minimizing
chart the pulled-back ideal is generated by $f_j$, of value $v(I)$.
\end{proof}

For a prime divisor $E$ over the closed point on a normal birational
model of a component of $\operatorname{Spec}R$, the change of primary boundary sends
\[
             \frac{\ord_E}{\ord_E(\mathfrak a)}
                    \longmapsto
             \frac{\ord_E}{\ord_E(\mathfrak b)}.
\]
Both denominators are positive and finite. On a chart where the two
pulled-back ideals are principal monomials
$u\prod_i z_i^{a_i}$ and $u'\prod_i z_i^{b_i}$, with $u,u'$ units,
a monomial valuation of weights $w_i$ satisfies
\[
                  v(\mathfrak a)=\sum_i a_iw_i,
                  \qquad v(\mathfrak b)=\sum_i b_iw_i.
\tag{70.7d}\label{eq:p121-monomial-gauge-slices}
\]
The two slices are the corresponding level-one hyperplanes; their change
rescales weights by $(\sum_i b_iw_i)^{-1}$. These coordinate formulas
require the named ideals to be monomial on the chosen chart. They do not
claim a resolution in arbitrary characteristic. Without that hypothesis,
the intrinsic ideal-value formula still applies, but coefficients along
the chart boundary alone need not compute a nonmonomial ideal's value.
For instance, in $k[[x,y]]$ the primary ideal $(x+y,x^2)$ has order
zero along both coordinate divisors, but value one at the monomial
valuation of weights $(1,1)$.

For a proper nonzero ideal $\mathfrak d$, the exact domain of its
normalization is the open scaling-invariant subcone
\[
 \mathcal C_R(\mathfrak d)
       =\{v\in\mathcal C_R:v(\mathfrak d)<+\infty\}.
\tag{70.7e}\label{eq:p121-nonprimary-gauge-domain}
\]
Indeed, $v(\mathfrak d)\geq v(\mathfrak m)>0$, and its finite-value
locus is open. Rescaling gives a homeomorphism from the corresponding
open subset of any primary section onto
$\mathcal V_{\mathfrak d}(R)=\{v\in\mathcal C_R:v(\mathfrak d)=1\}$.
Its inverse rescales by the primary ideal value. For arbitrary ideals
$\mathfrak d,\mathfrak e$, the same formula defines a partial change
on $\{v\in\mathcal V_{\mathfrak d}:0<v(\mathfrak e)<+\infty\}$,
with image $\{u\in\mathcal V_{\mathfrak e}:0<u(\mathfrak d)<+\infty\}$.
The inverse and cocycle identities hold on these domains and commute
with normalization-branch restriction. The unit ideal has zero gauge,
the zero ideal has infinite gauge, and neither has a level-one section.

\begin{example}\label{ex:p121-nonprimary-noncompact-link}
For $R=k[[x,y]]$ and $\mathfrak d=(x)$, the semivaluation from
$R/(x)=k[[y]]$ and $\ord_y$ lies in $\mathcal V_{\mathfrak m}(R)$
but has infinite $\mathfrak d$-value. Its ray is omitted by
\eqref{eq:p121-nonprimary-gauge-domain}. The monomial valuations
$v_j(x)=1$, $v_j(y)=1/j$ belong to $\mathcal V_{\mathfrak d}(R)$.
If this link were compact, its image under continuous $y$-evaluation
would be a closed compact subset of $[0,+\infty]$ containing every
$1/j$, hence zero. But every centered semivaluation is strictly positive
on $y\in\mathfrak m$. Thus this link is not compact.
\end{example}

Finally let $R$ be nonreduced, with nilradical $\mathcal N$.
If $n^k=0$, then $kv(n)=+\infty$, so $v(n)=+\infty$.
Applying the ultrametric inequality to $f+n$ and to $(f+n)-n$
gives $v(f+n)=v(f)$. Consequently reduction induces a homeomorphism
of centered cones and preserves every ideal gauge and change of section.
If $\dim R_{\mathrm{red}}=0$, that reduced local ring is a field,
so its centered cone is empty. Nilpotent multiplicities cannot be
recovered from this datum and must be retained separately for any
intersection or measure construction.

These results settle the normalizing-boundary question for excellent
local germs: the ray space, primary sections, partial nonprimary
sections, and conductor descent are explicit. They do not construct
a singular psh class, envelopes, products, energy, or general finite-map
pluripotential functoriality.
